\begin{helper}
Let \(L\) and \(\alpha\) be real parameters.  Suppose the one-dimensional
exponential tail calculation gives the numeric lower bound
\[
\frac53+2\alpha
\le
2(1-\alpha)^2 e^{-\alpha L}
\int_0^L \bigl(e^{-x}-e^{-L}\bigr)\,dx .
\]
Then the high-overlap exponential window has the required margin:
\[
\frac53+2\alpha
\le
2\int_0^L\int_x^L
(1-\alpha)e^{-\alpha L}(1-\alpha)e^{-y}\,dy\,dx .
\]
\end{helper}
\begin{proof}
For fixed \(x\), the only \(y\)-dependent factor in the inner integral is
\(e^{-y}\).  The derivative of \(-e^{-y}\) is \(e^{-y}\), so the fundamental
theorem of calculus gives
\[
\int_x^L e^{-y}\,dy=e^{-x}-e^{-L}.
\]
Pulling the \(y\)-independent factors out of the inner integral therefore gives
\[
\int_x^L
(1-\alpha)e^{-\alpha L}(1-\alpha)e^{-y}\,dy
=
(1-\alpha)^2 e^{-\alpha L}\bigl(e^{-x}-e^{-L}\bigr).
\]
Integrating this identity in \(x\) from \(0\) to \(L\), and pulling the constant
\((1-\alpha)^2e^{-\alpha L}\) outside the outer integral, yields
\[
\int_0^L\int_x^L
(1-\alpha)e^{-\alpha L}(1-\alpha)e^{-y}\,dy\,dx
=
(1-\alpha)^2 e^{-\alpha L}
\int_0^L \bigl(e^{-x}-e^{-L}\bigr)\,dx .
\]
Multiplying by \(2\) and applying the assumed one-dimensional lower bound gives
the displayed high-overlap margin.
\end{proof}
